\documentclass[11pt]{article}

% Change "review" to "final" to generate the final (sometimes called camera-ready) version.
% Change to "preprint" to generate a non-anonymous version with page numbers.
\usepackage[review]{acl}

% Standard package includes
\usepackage{times}
\usepackage{latexsym}
\usepackage{graphicx}
\usepackage{amsmath}
\usepackage{booktabs}
\usepackage{float}
\usepackage[T1]{fontenc}
\usepackage[utf8]{inputenc}
\usepackage{microtype}
\usepackage{inconsolata}
\raggedbottom

% If the title and author information does not fit in the area allocated, uncomment the following
%
%\setlength\titlebox{<dim>}
%
% and set <dim> to something 5cm or larger.

\title{The Token Tax: Evaluating Script Fairness and Tokenizer Fertility in Telugu}

% Author information can be set in various styles:
% For several authors from the same institution:
% \author{Author 1 \and ... \and Author n \\
%         Address line \\ ... \\ Address line}
% if the names do not fit well on one line use
%         Author 1 \\ {\bf Author 2} \\ ... \\ {\bf Author n} \\
% For authors from different institutions:
% \author{Author 1 \\ Address line \\  ... \\ Address line
%         \And  ... \And
%         Author n \\ Address line \\ ... \\ Address line}
% To start a separate ``row'' of authors use \AND, as in
% \author{Author 1 \\ Address line \\  ... \\ Address line
%         \AND
%         Author 2 \\ Address line \\ ... \\ Address line \And
%         Author 3 \\ Address line \\ ... \\ Address line}

\author{First Author \\
  Affiliation / Address line 1 \\
  Affiliation / Address line 2 \\
  Affiliation / Address line 3 \\
  \texttt{email@domain} \\\And
  Second Author \\
  Affiliation / Address line 1 \\
  Affiliation / Address line 2 \\
  Affiliation / Address line 3 \\
  \texttt{email@domain} \\}

%\author{
%  \textbf{First Author\textsuperscript{1}},
%  \textbf{Second Author\textsuperscript{1,2}},
%  \textbf{Third T. Author\textsuperscript{1}},
%  \textbf{Fourth Author\textsuperscript{1}},
%\\
%  \textbf{Fifth Author\textsuperscript{1,2}},
%  \textbf{Sixth Author\textsuperscript{1}},
%  \textbf{Seventh Author\textsuperscript{1}},
%  \textbf{Eighth Author \textsuperscript{1,2,3,4}},
%\\
%  \textbf{Ninth Author\textsuperscript{1}},
%  \textbf{Tenth Author\textsuperscript{1}},
%  \textbf{Eleventh E. Author\textsuperscript{1,2,3,4,5}},
%  \textbf{Twelfth Author\textsuperscript{1}},
%\\
%  \textbf{Thirteenth Author\textsuperscript{3}},
%  \textbf{Fourteenth F. Author\textsuperscript{2,4}},
%  \textbf{Fifteenth Author\textsuperscript{1}},
%  \textbf{Sixteenth Author\textsuperscript{1}},
%\\
%  \textbf{Seventeenth S. Author\textsuperscript{4,5}},
%  \textbf{Eighteenth Author\textsuperscript{3,4}},
%  \textbf{Nineteenth N. Author\textsuperscript{2,5}},
%  \textbf{Twentieth Author\textsuperscript{1}}
%\\
%\\
%  \textsuperscript{1}Affiliation 1,
%  \textsuperscript{2}Affiliation 2,
%  \textsuperscript{3}Affiliation 3,
%  \textsuperscript{4}Affiliation 4,
%  \textsuperscript{5}Affiliation 5
%\\
%  \small{
%    \textbf{Correspondence:} \href{mailto:email@domain}{email@domain}
%  }
%}

\begin{document}
\maketitle
\begin{abstract}
Telugu has more than 96 million speakers but remains low resource in
LLM development, underrepresented in pretraining data and absent
from tokenizer design. This creates a \textit{token tax}: native
script users pay more per API call and get a smaller context window
than English speakers for the same content. I audit six tokenizers
across three corpus registers, native formal, native informal, and
romanized informal Telugu, 2,789 sentences, plus 62 native speaker
validated minimal pairs across nine agglutinative morphological
categories. Claude Haiku~4.5 needs 7.73 tokens per word on native
Telugu; Anon-Indic-GPT2 needs only 1.60, a 4.8$\times$ gap from
training data. The legacy GPT-2 tokenizer costs 4.66$\times$ more for
native script than romanized text (Wilcoxon signed rank,
$p < 10^{-165}$, 100\% of pairs); Claude Haiku~4.5 carries a smaller
1.56$\times$ penalty (99.9\%). Models trained on Indic data reverse
this bias. At current prices, Telugu speakers pay roughly
7.7$\times$ more than English speakers for equivalent text, with a
context window 7.7$\times$ smaller. On MILU
\citep{verma-etal-2024-milu} (7,304 Telugu MCQ questions), efficiency
and accuracy are decoupled: GPT-4o scores 72.53\% (5 shot) despite a
moderate fertility penalty, while Anon-Indic-GPT2, the most
efficient tokenizer, scores 26.97\%, near chance. The token tax
reflects pricing, not capability.
\end{abstract}

\section{Introduction}
When an English speaker and a Telugu speaker send semantically
equivalent text to a commercial API, they end up paying different
amounts, not because of what the text says, but because of which
script it is written in. Tokenizers, the components that turn raw
text into the integers a model actually reads, are trained mostly on
English and other high resource European languages. Because of this,
they routinely over-segment the morphologically rich scripts of low
resource languages, producing far more tokens per word. That
difference is the hidden \textit{token tax} those users pay.

Telugu makes this problem concrete. Spoken by over 96 million people
across India and its diaspora \citep{census2011}, Telugu is a
Dravidian language with agglutinative morphology: a single word can
encode tense, aspect, number, case, and honorificity all at once
through chains of suffixes. These are exactly the kind of surface
forms that general purpose tokenizers struggle to segment well.
Despite its large speaker base, Telugu is still low resource in the
NLP sense: it is underrepresented in pretraining corpora, largely
absent from tokenizer design decisions, and rarely included in
multilingual benchmarks except as a single row in a large table.

Prior work has shown that tokenizer fertility (tokens per word)
varies a lot across languages \citep{rust-etal-2021-how}, and that
this variation carries real consequences for task performance, API
cost, and how much of the context window a user actually gets to use
\citep{petrov-etal-2023-language, ahia-etal-2023-languages}. But
existing work tends to treat Telugu only in passing, lumps native
script and romanized Telugu together as one variety, and never
breaks fertility penalties down by morphological type. The script
dimension of fairness, whether a tokenizer charges people more for
writing in their own script than for romanizing it, has not really
been studied on its own terms.

This paper is a focused empirical audit aimed at that gap. I
evaluate six tokenizers, general purpose commercial APIs, open
source baselines built for English, and tokenizers built specifically
for Indic languages, across three corpus registers and a
morphologically annotated minimal pair dataset. Four questions
guide the analysis:
\begin{itemize}
    \item \textbf{RQ1} (Fertility): How much does tokenizer fertility
    vary across tokenizers and corpus registers for Telugu?
    \item \textbf{RQ2} (Script fairness): Is there a statistically
    significant gap between native script and romanized fertility,
    and which direction does it run?
    \item \textbf{RQ3} (Morphology): Which agglutinative morphological
    structures incur the greatest token-count penalty, and does the
    penalty vary systematically across tokenizer tiers?
    \item \textbf{RQ4} (Decoupling): Does lower fertility predict
    higher downstream task accuracy, or are tokenizer efficiency and
    model capability independent of each other?
\end{itemize}
My contributions are:
\begin{itemize}
    \item A fertility audit of six tokenizers across 2,789 sentences
    in three Telugu registers, including per sentence Wilcoxon
    signed rank tests and bootstrap confidence intervals on script
    fairness gaps (\S\ref{sec:results-fertility}, \S\ref{sec:results-gap});
    \item A morphological breakdown built from 62 minimal pairs
    validated by native speakers, spanning nine agglutinative
    categories, that identifies which linguistic structures take the
    biggest penalty (\S\ref{sec:results-minimal});
    \item A zero shot MILU evaluation of Claude Haiku~4.5 across
    7,304 Telugu MCQ questions, showing that tokenizer efficiency and
    downstream accuracy are decoupled (\S\ref{sec:results-milu});
    \item A quantification of the real world cost and context
    window consequences Telugu speakers face when using commercial
    APIs (\S\ref{sec:discussion}).
\end{itemize}
I find that tokenizers trained on Indic data not only reach lower
fertility but actually reverse the script fairness gap, making
native Telugu cheaper to tokenize than romanized input. General
purpose tokenizers do the opposite, imposing statistically
significant penalties as high as 4.66$\times$. These results matter
directly for how fairly speakers of morphologically rich, low
resource languages can access AI systems.



\section{Related Work}

\paragraph{Tokenizer fertility and cross lingual fairness.}
\citet{rust-etal-2021-how} were the first to systematically measure
tokenizer fertility across languages. They showed that multilingual
models underperform on languages with high tokens per word ratios
because of representational fragmentation. \citet{petrov-etal-2023-language}
extended this analysis to commercial tokenizers. They found that
speakers of low resource languages are charged more tokens, and
therefore more money, for semantically equivalent content: a
structural unfairness built into API pricing.
\citet{ahia-etal-2023-languages} quantified this disparity across
164 languages and found that some languages cost up to 10$\times$
more to process than English. My work narrows this line of inquiry
to Telugu specifically and adds a \textit{script dimension} (native
versus romanized) that prior cross lingual surveys do not address.

\paragraph{Tokenization for morphologically rich languages.}
Agglutinative languages pose a fundamental challenge for BPE based
tokenizers: suffix chains that encode grammatical information as
single morphemes get split into multiple subword units. This
destroys morphological coherence \citep{bostrom-durrett-2020-byte}.
\citet{bostrom-durrett-2020-byte} further show that BPE
systematically fails on agglutinative languages by splitting across
word boundaries, with the resulting token sequences obscuring
grammatical structure.
\citet{gutierrez-vasques-etal-2021-characters} trace this failure to
the merge order used in BPE training, which favors high frequency
English subwords. These findings motivate my own morphological
breakdown using minimal pairs, which isolates the agglutinative
structures that incur the greatest fertility penalty.

\paragraph{Indic NLP and Telugu resources.}
The AI4Bharat initiative has produced the primary infrastructure for
Indic language NLP, including IndicBERT
\citep{kakwani-etal-2020-indicnlpsuite}, IndicXTREME
\citep{doddapaneni-etal-2023-towards}, and the MILU benchmark
\citep{verma-etal-2024-milu} used in my evaluation.
\citet{kunchukuttan-2020-indicnlp} provides the transliteration
tools behind my romanized corpus. Even with this growing ecosystem,
Telugu tokenization has not been studied on its own: it shows up as
one row among 12 to 22 Indic languages in aggregate analyses, with
no breakdown by register or script. The Sarvam family of models
\citep{sarvam-2024} is the first commercially deployed tokenizer
built specifically for Indic languages at scale, which makes it a
natural point of comparison. My audit is the first to evaluate and
compare all three tiers, general purpose commercial, English open
source, and Indic specialized, for Telugu in one controlled
experimental setting.

\section{Methodology}
\label{sec:methodology}

\subsection{Corpus Construction}

I build a three register Telugu corpus that supports both aggregate
fertility analysis and paired significance testing across script
varieties.

\paragraph{Native formal.}
I sample 789 sentences from the Telugu Wikipedia corpus
\texttt{vengi-ai/telugu-wikipedia-clean} \citep{vengi-ai-2023},
keeping only sentences with 5 to 50 words separated by white space.
Wikipedia prose is a good stand in for formal written Telugu: dense
vocabulary, almost no romanization.

\paragraph{Native informal.}
I sample 1,000 sentences from Telugu sentiment analysis and hate
speech detection datasets on the HuggingFace Hub, keeping only
entries written in native script. Informal registers show more
morphological variation, colloquial contractions, and code
switching, none of which show up much in formal prose.

\paragraph{Romanized informal.}
I build a line by line romanized version of the native informal
corpus using systematic ITRANS transliteration via the
\texttt{indic-transliteration} library
\citep{kunchukuttan-2020-indicnlp}. Each romanized sentence lines up
exactly with its native script counterpart, which lets me run
paired significance tests sentence by sentence. ITRANS is the
romanization scheme most commonly used for informal Telugu online.

\paragraph{Preprocessing.}
Every register goes through the same cleaning pipeline: regex based
removal of personally identifiable information, punctuation stripping
to prevent artificial token inflation, deduplication on normalized
surface form, and length filtering. The final corpus has
2,789 sentences in total: 789 native formal, 1,000 native informal,
and 1,000 romanized informal.

\subsection{Tokenizers}

I evaluate six tokenizers across three tiers of Indic language
coverage; Table~\ref{tab:tokenizers} summarizes them.
\begin{table}[H]
\centering
\resizebox{\columnwidth}{!}{%
\begin{tabular}{llr}
\hline
\textbf{Tokenizer} & \textbf{Backend} & \textbf{Vocab} \\
\hline
GPT-4o                & tiktoken \texttt{o200k\_base}  & 200k \\
Claude Haiku~4.5      & Anthropic API                  & n/d\textsuperscript{$\dagger$}  \\
GPT-2 (legacy)        & HuggingFace BPE                & 50k  \\
Sarvam-2B             & HuggingFace SentencePiece      & 64k  \\
Sarvam-105B           & HuggingFace SentencePiece      & 64k  \\
Anon-Indic-GPT2 & SentencePiece                  & 32k  \\
\hline
\end{tabular}%
}
\caption{Tokenizers, backends, and vocabulary sizes.
         \textsuperscript{$\dagger$}Not publicly disclosed by Anthropic.}
\label{tab:tokenizers}
\end{table}
\paragraph{GPT-4o.}
I get token counts from the \texttt{tiktoken} library using the
\texttt{o200k\_base} encoding, the same tokenizer that GPT-4o uses.

\paragraph{Claude Haiku 4.5.}
I get token counts through the Anthropic
\texttt{messages.count\_tokens} API, using the model identifier
\texttt{claude-haiku-4-5-20251001} (Claude Haiku~4.5, released
October 2025). This is a separate model from Claude~3.5 Haiku,
Anthropic's fourth generation Haiku release. For the MILU benchmark
I use the same model through the Anthropic Messages API, accessed
in June 2026.

\paragraph{GPT-2 (legacy).}
I load the original GPT-2 tokenizer \citep{radford-etal-2019-gpt2}
through \texttt{transformers.AutoTokenizer} from the \texttt{gpt2}
HuggingFace repository. This BPE tokenizer was built around English
and serves as my worst case baseline. It encodes Telugu Unicode
characters at the byte level, which drives its fertility extremely
high.

\paragraph{Sarvam-2B and Sarvam-105B.}
I load both Sarvam models \citep{sarvam-2024} via
\texttt{AutoTokenizer.from\_pretrained} using
\texttt{sarvamai/sarvam-1-v0.5} and \texttt{sarvamai/sarvam-105b}
respectively. Both were trained on large Indic corpora and use a
SentencePiece vocabulary of 64k that covers all 22 scheduled Indian
languages.

\paragraph{Anon-Indic-GPT2.}
This model \citep{anonymous-2026-indic} is a Telugu-specialized
GPT-2 variant with a compact BPE vocabulary trained on Telugu text.
I load the tokenizer via its custom vocabulary interface.

\subsection{Fertility and Script Fairness Metrics}

\paragraph{Fertility.}
Following \citet{rust-etal-2021-how}, I define the fertility of
tokenizer $t$ on register $R$ as:

\begin{equation}
    \mathcal{F}(t, R) =
        \frac{\displaystyle\sum_{s \in R} \mathrm{tokens}(t, s)}
             {\displaystyle\sum_{s \in R} \mathrm{words}(s)}
    \label{eq:fertility}
\end{equation}

Here $\mathrm{words}(s)$ is the word count of sentence $s$, counted
by splitting on white space. Lower fertility means more efficient
tokenization.

\paragraph{Script fairness gap.}
I define the script fairness gap for tokenizer $t$ as the ratio of
native script fertility to romanized fertility on the matched
registers:

\begin{equation}
    \mathcal{G}(t) =
        \frac{\mathcal{F}(t,\ \mathrm{native\_informal})}
             {\mathcal{F}(t,\ \mathrm{romanized\_informal})}
    \label{eq:gap}
\end{equation}

A value of $\mathcal{G}(t) > 1$ means native script costs more to
tokenize than romanized input; a value below 1 means the opposite.
Because the romanized corpus is a line by line transliteration of
the native corpus, both registers carry the same meaning. That
makes this ratio a direct measure of script level tokenizer bias.

\paragraph{Significance testing.}
To test whether the script fairness gap is statistically
significant, I run a one-sided Wilcoxon signed rank test
\citep{wilcoxon-1945} ($H_1$: native fertility $>$ romanized
fertility) on the per sentence fertility gaps
$\delta_i = \mathcal{F}(t, s_i^{\mathrm{nat}})
- \mathcal{F}(t, s_i^{\mathrm{rom}})$ across the 1,000 matched
sentence pairs ($n = 1{,}000$). I also compute 95\% bootstrap
confidence intervals on the median gap, using 1,000 resamples with
replacement.

\subsection{Minimal Pair Analysis}

To find out which morphological structures take the biggest
fertility hit, I build a minimal pair dataset of 62 Telugu word
forms validated by native speakers, spanning nine agglutinative
categories: \textit{base\_noun}, \textit{case\_suffix},
\textit{case\_suffix\_with\_sandhi}, \textit{plural\_suffix},
\textit{honorific\_suffix}, \textit{compound\_with\_sandhi},
\textit{borrowed\_base}, \textit{borrowed\_plus\_case\_suffix}, and
\textit{verb\_agglutination\_chain}. I tokenize each word with all
six tokenizers and report the mean and median token count per
category. I release the dataset with my repository so others can
build on it.

\subsection{Benchmark Evaluation}
\label{sec:benchmark}

I evaluate downstream Telugu understanding on MILU
\citep{verma-etal-2024-milu}, a 7,304 question multiple choice
benchmark that covers 41 subjects across 8 domains in the Telugu
test split. I use three evaluation protocols, each matched to what
the model can actually do.

\paragraph{Zero shot generative (Claude Haiku 4.5).}
Each question is formatted as a four option MCQ prompt. I use the
Anthropic assistant prefill technique: I append
\texttt{\small\{"role": "assistant", "content": "The answer is"\}}
to the message list to push the model toward a single letter
response, with \texttt{max\_tokens=50} and greedy decoding. I parse
responses with a multi strategy letter extractor; 7,108 of the
7,304 responses (97.3\%) give me a parseable letter. The 196
unparseable responses (2.7\%) are excluded from the accuracy
calculation and show no systematic concentration in any subject
domain.

\paragraph{Log likelihood (Anon-Indic-GPT2).}
Since Anon-Indic-GPT2 is a base language model without
instruction following ability, I evaluate it with the log
likelihood method: for each option $o_k$, I compute the length
normalised mean token log probability of the option's tokens given
the question prefix, then pick the option with the highest score:

\begin{equation}
    \hat{k} = \mathop{\arg\max}_{k}\
              \frac{1}{|o_k|}
              \sum_{j=1}^{|o_k|}
              \log P\!\left(o_{k,j} \mid q,\, o_{k,<j}\right)
    \label{eq:loglik}
\end{equation}

\paragraph{GPT-4o (paper score).}
I report GPT-4o's 5 shot accuracy from Table~3 of
\citet{verma-etal-2024-milu} (72.53\%). Running my own full
inference pass was too expensive, so I flag the methodological
difference (5 shot versus zero shot) wherever it matters.

\paragraph{Sarvam-2B (excluded).}
Zero shot generative evaluation of Sarvam-2B yielded only 576
parseable responses out of 7,304 (7.9\%). The model answered in
Telugu prose instead of giving a letter, which fits with the fact
that it was not instruction tuned on English language MCQ formats.
I exclude Sarvam-2B from my accuracy analysis as a result, and
cite the MILU paper's 5 shot score of 28.57\% for reference only.


\section{Results}
\label{sec:results}

\subsection{RQ1: Tokenizer Fertility Across Registers}
\label{sec:results-fertility}

Table~\ref{tab:fertility} shows fertility $\mathcal{F}(t, R)$ for
each of the six tokenizers across the three corpus registers.
\begin{table}[H]
\centering
\resizebox{\columnwidth}{!}{%
\begin{tabular}{@{}lrrr@{}}
\toprule
\textbf{Tokenizer} & \textbf{Native Formal} &
\textbf{Native Informal} & \textbf{Romanized Informal} \\
\midrule
GPT-4o                &  3.28 &  3.25 &  3.98 \\
Claude Haiku~4.5      &  6.91 &  7.73 &  4.95 \\
GPT-2 (legacy)        & 19.99 & 22.08 &  4.74 \\
Sarvam-2B             &  2.52 &  3.01 &  5.67 \\
Sarvam-105B           &  2.68 &  2.71 &  4.01 \\
Anon-Indic-GPT2 & \textbf{1.68} & \textbf{1.60} &  5.49 \\
\bottomrule
\end{tabular}%
}
\caption{Fertility (tokens/word) by tokenizer and register. Bold = best.}
\label{tab:fertility}
\end{table}
Fertility varies a lot across tokenizers. On the native informal
register, which is the closest proxy I have for everyday Telugu
usage, Anon-Indic-GPT2 reaches the lowest fertility at 1.60
tokens per word, while GPT-2 reaches 22.08. That is a 13.8$\times$
gap between the most and least efficient tokenizers on the same
text. Claude Haiku~4.5 produces 7.73 tokens per word, 4.8$\times$
higher than Anon-Indic-GPT2.

The Sarvam models, both trained on Indic data, reach competitive
fertility (2.71 to 3.01) on native informal text, close to GPT-4o's
3.25, even though they target a broader vocabulary covering 22
Indic languages. This suggests that deliberate multilingual
tokenizer design can close most of the fertility gap without giving
up cross lingual coverage.

\paragraph{Token length distribution.}
The token length distribution explains where this gap comes from
structurally. Claude Haiku~4.5 assigns five or more tokens to 82.5\%
of native informal Telugu words. The legacy GPT-2 tokenizer assigns
five or more tokens to 95.7\% of words; at that rate it is
essentially falling back to byte level encoding for Telugu Unicode.
Anon-Indic-GPT2, by contrast, keeps 86.2\% of words within one
or two tokens (61.6\% as a single token). That fits with segmentation close to the morpheme
level, in a language where words are typically just one to three
morphemes long.


\subsection{RQ2: Script Fairness Gap}
\label{sec:results-gap}

Table~\ref{tab:gap} reports the script fairness gap $\mathcal{G}(t)$
together with Wilcoxon signed rank test results and 95\% bootstrap
confidence intervals on the median per sentence gap.
\begin{table}[H]
\centering
\resizebox{\columnwidth}{!}{%
\begin{tabular}{@{}lrrrr@{}}
\toprule
\textbf{Tokenizer} & $\mathcal{G}(t)$ &
\textbf{Med.\ $\delta$} & \textbf{95\% CI} & \textbf{$p$} \\
\midrule
GPT-4o                & 0.82 & $-$0.73 & [$-$0.76,\;$-$0.69] & 1.0 \\
Claude Haiku~4.5      & 1.56 & $+$2.75 & [$+$2.71,\;$+$2.81] & $<\!10^{-165}$ \\
GPT-2 (legacy)        & \textbf{4.66} & $+$17.22 & [$+$17.0,\;$+$17.4] & $<\!10^{-165}$ \\
Sarvam-2B             & 0.53 & $-$2.65 & [$-$2.70,\;$-$2.58] & 1.0 \\
Sarvam-105B           & 0.68 & $-$1.30 & [$-$1.33,\;$-$1.25] & 1.0 \\
Anon-Indic-GPT2 & \textbf{0.29} & $-$3.88 & [$-$4.00,\;$-$3.83] & 1.0 \\
\bottomrule
\end{tabular}%
}
\caption{Script fairness gap $\mathcal{G}(t)$, median $\delta$,
         95\% CI, and Wilcoxon $p$ ($n=1{,}000$).
         $\mathcal{G}(t){>}1$: native costlier;
         $\mathcal{G}(t){<}1$: romanized costlier.}
\label{tab:gap}
\end{table}
The results split cleanly along tokenizer origin. Western, general
purpose tokenizers penalize native Telugu script: the legacy GPT-2
tokenizer imposes a 4.66$\times$ penalty, with native script
costlier in 100\% of matched pairs ($p < 10^{-165}$), and Claude
Haiku~4.5 imposes a 1.56$\times$ penalty (99.9\% of pairs,
$p < 10^{-165}$). On these tokenizers, a Telugu speaker who writes
in their own script pays substantially more tokens than one who
romanizes the same content, which means more money and more latency
for no difference in meaning.

Tokenizers trained on Indic data flip this bias the other way
entirely. Sarvam-2B ($\mathcal{G} = 0.53$), Sarvam-105B
($\mathcal{G} = 0.68$), and Anon-Indic-GPT2 ($\mathcal{G} =
0.29$) all make native script cheaper than romanized input, with
native script costlier in only 0.0\% to 0.2\% of pairs.
Anon-Indic-GPT2 shows the strongest effect: native Telugu
script costs 3.4$\times$ less than its romanized equivalent, the
largest native script advantage I observed. GPT-4o sits in the
middle ($\mathcal{G} = 0.82$), with a small romanized advantage that
is not statistically significant ($p=1.0$).

\subsection{RQ3: Morphological Breakdown}
\label{sec:results-minimal}

Table~\ref{tab:minimal} gives mean tokens per word for each
morphological category, across all six tokenizers.
\begin{table}[H]
\centering
\resizebox{\columnwidth}{!}{%
\begin{tabular}{@{}lrrrrrr@{}}
\toprule
\textbf{Morphological Type} & \textbf{G4o} & \textbf{Cla.} &
\textbf{G2} & \textbf{S-2B} & \textbf{S-105B} & \textbf{Anon.} \\
\midrule
base\_noun                    & 2.6 &  4.4 & 12.4 & 2.1 & 2.0 & \textbf{1.0} \\
borrowed\_base                & 3.5 &  6.3 & 17.0 & 2.0 & 2.8 & \textbf{1.2} \\
borrowed\_plus\_case\_suffix  & 3.7 &  8.8 & 25.5 & 2.8 & 3.5 & \textbf{1.7} \\
case\_suffix                  & 3.8 &  5.8 & 19.2 & 3.0 & 3.2 & \textbf{1.6} \\
case\_suffix\_with\_sandhi    & 3.7 &  7.3 & 23.4 & 2.5 & 3.0 & \textbf{1.7} \\
compound\_with\_sandhi        & 4.8 &  9.0 & 27.0 & 3.3 & 4.2 & \textbf{1.8} \\
honorific\_suffix             & 5.7 & 11.7 & 35.5 & 2.8 & 4.5 & \textbf{2.2} \\
plural\_suffix                & 3.4 &  7.1 & 21.0 & 2.0 & 3.1 & \textbf{1.3} \\
verb\_agglutination\_chain    & 4.6 &  8.8 & 29.6 & 2.0 & 3.5 & \textbf{1.8} \\
\bottomrule
\end{tabular}%
}
\caption{Mean tokens/word by morphological type.
         G4o=GPT-4o, Cla.=Claude Haiku~4.5, G2=GPT-2,
         S-2B=Sarvam-2B, S-105B=Sarvam-105B, Anon.=Anon-Indic-GPT2.
         Bold = lowest per row.}
\label{tab:minimal}
\end{table}
Two morphological categories take the worst hit. \textit{Honorific
suffix} forms, where an honorific morpheme attaches to a noun or
verb stem, average 35.5 tokens per word under GPT-2 and 11.7 under
Claude Haiku~4.5, the highest counts anywhere in the study.
\textit{Verb agglutination chains}, which stack tense, aspect, mood,
and agreement morphemes onto a single verb, average 29.6 under GPT-2
and 8.8 under Claude. Both are exactly the kind of construction that
defines Telugu morphology: common in natural text, and carrying more
grammatical information per surface word than almost anything else
in the language.

Sarvam-2B handles every morphological category in 2.0 to 3.3 tokens
per word, holding its own against or beating GPT-4o despite covering
22 languages instead of just one. Anon-Indic-GPT2 does even
better, staying in the 1.0 to 2.2 range across every category, close
to true morpheme level segmentation throughout. It tokenizes every
\textit{base\_noun} form in my set as a single token (mean = 1.0).
That points to a vocabulary large enough to store common Telugu
stems directly, rather than splitting them.

\subsection{RQ4: Benchmark Accuracy and the Decoupling Finding}
\label{sec:results-milu}

Table~\ref{tab:milu} reports MILU Telugu test accuracy for the
models I have scores for. Sarvam-2B is excluded here because its
7.9\% parse rate (576/7,304) makes its generative accuracy
uninterpretable; I do not report the MILU paper's 5-shot score for
it since the evaluation setup is not comparable to my runs.
\begin{table}[H]
\centering
\resizebox{\columnwidth}{!}{%
\begin{tabular}{@{}llrr@{}}
\hline
\textbf{Model} & \textbf{Method} &
\textbf{Fertility} & \textbf{Accuracy} \\
\hline
GPT-4o                & 5-shot (paper)      & 3.25 & 72.53\% \\
Claude Haiku~4.5      & 0-shot (ours)       & 7.73 & 60.78\% \\
Anon-Indic-GPT2 & Log-likelihood (ours) & 1.60 & 26.97\% \\
\hline
Chance baseline       & ---                 & ---  & 25.00\% \\
\hline
\end{tabular}%
}
\caption{MILU Telugu accuracy and native informal fertility.
         Evaluation protocols differ across models (see text).}
\label{tab:milu}
\end{table}
These numbers are consistent with tokenizer fertility and downstream
accuracy being decoupled, though the comparison is limited by
incompatible evaluation protocols: GPT-4o's score is 5-shot from
\citet{verma-etal-2024-milu}, while mine are zero-shot generative
and log-likelihood. GPT-4o gets the highest accuracy, 72.53\%, with only
a moderate fertility of 3.25. Claude Haiku~4.5 scores 60.78\%
despite having the worst fertility among the instruction tuned
models I tested (7.73). Anon-Indic-GPT2 has the best
fertility in the study (1.60) and still scores near chance, 26.97\%
against a 25.0\% baseline. That is not a tokenizer problem: it is a
110M parameter base language model with no instruction tuning, too
small for the task regardless of how cheaply it tokenizes.

The pattern has a simple explanation. Tokenizer efficiency and model
capability come from separate design choices, so there is no reason
to expect one to track the other. A model can be trained on large
amounts of Indic data and end up with an efficient tokenizer while
still being a small base model. A frontier model can do the
opposite: understand Telugu well while carrying a poorly calibrated
tokenizer. I read the token tax as a pricing problem sitting at the
tokenizer level, separate from whatever the underlying model can
do, and it needs to be fixed there rather than waiting on bigger
models.

\paragraph{Cost and context implications.}
At Claude Haiku~4.5's pricing of \$0.0008 per 1,000 input tokens, a
fertility of 7.73 tokens per word on native Telugu means processing
1,000 Telugu words costs about \$0.0062, roughly 7.7$\times$ more
than the same amount of English text at close to 1.0 token per
word.\footnote{English fertility of $\approx$1.0 tokens/word is a
standard approximation from prior work
\citep{petrov-etal-2023-language}; I do not measure it directly in
this study.} The 200,000 token context window works out to about 25,873
native Telugu words, against roughly 200,000 English words: the
effective window shrinks by the same 7.7$\times$. GPT-4o's fertility
of 3.25 translates to a 3.25$\times$ penalty in both cost and
context. None of this is incidental; it adds another layer of
disadvantage on top of everything else that comes with speaking a
low resource language with little representation in model training
data.




\section{Discussion}
\label{sec:discussion}

\subsection{The Token Tax as a Pricing Equity Problem}

My results reframe tokenizer fertility as a question of economic
access, not model capability. The term \textit{token tax} captures
this directly: Telugu speakers who write in their native script are
charged more tokens for input that means exactly the same thing as a
sentence in a well represented language, and that extra token count
turns into more money, more latency, and less usable context.

How big this tax is depends entirely on which tokenizer the
speaker's application happens to use. At Claude Haiku~4.5 pricing, a
fertility of 7.73 tokens per word on native Telugu means a
7.7$\times$ cost premium over English, which sits close to 1.0
token per word. At GPT-4o's 3.25 tokens per word, the premium drops
to 3.25$\times$. These multipliers stack across every interaction: a
500 word Telugu message that costs \$0.003 in English costs \$0.024
on Claude Haiku~4.5, and the effective context window shrinks by the
same factor. For someone paying for AI access out of a prepaid plan,
which is how most people in India get online, this is not an
abstract number. It is a hard ceiling on how much of the model they
can actually use per rupee.

The token tax also lands on speakers who already face the biggest
barriers to AI access. Telugu has among the lowest per capita LLM
representation of any of India's 22 scheduled languages
\citep{doddapaneni-etal-2023-towards}. Telugu speakers therefore pay
more per token and get models trained on less of their own language
at the same time. Neither problem causes the other, but together
they compound.

\subsection{Tokenizer Design Drives the Gap}

The split between Western and Indic tokenizers in
Table~\ref{tab:gap} is not a coincidence. It comes directly from how
each tokenizer's vocabulary was built.

General purpose tokenizers, trained mostly on English and European
text, have no subword units built for Telugu morphemes. When they
hit a Telugu word, they fall back to character level or byte level
segmentation, and fertility goes up sharply along with the penalty
on native script. GPT-2 is the extreme case: its 50k vocabulary has
no Telugu specific tokens at all, so every Unicode character gets
encoded at the byte level. That alone produces the 4.66$\times$
script fairness gap. Claude Haiku~4.5's smaller 1.56$\times$ gap
tells a different story: its vocabulary does include some Telugu
subword units, which fits with it being trained on multiple
languages, but not enough of them to overcome a tokenizer built
around Latin script text.

Tokenizers trained on Indic data flip this relationship completely.
Sarvam-2B and Sarvam-105B land at gaps of 0.53 and 0.68, making
native script cheaper than romanized input. That makes sense
linguistically: native Telugu script packs more meaning into each
character than ITRANS romanization does, since ITRANS spells out
aspirated consonants and long vowels as several characters each. A
vocabulary trained on native text will naturally represent that text
more efficiently. Anon-Indic-GPT2 pushes this furthest, with a
gap of 0.29, because its entire 32k vocabulary was trained only on
Telugu text. Every slot in that vocabulary serves Telugu morphology;
none of it goes to other languages.

GPT-4o sits close to neutral on this measure ($\mathcal{G} = 0.82$),
which suggests a large enough multilingual vocabulary, 200k tokens
in this case, can get close to script fairness without being built
around any one language. It just needs enough Telugu subword units
in there to avoid falling back to bytes. That is a useful result: it
means tokenizer fairness does not require giving up cross lingual
coverage, as long as the vocabulary budget is big enough.

\subsection{Decoupling Fertility from Capability}

The MILU results in Table~\ref{tab:milu} show that tokenizer
efficiency and downstream accuracy do not move together. GPT-4o
gets the highest accuracy, 72.53\%, at a moderate fertility of 3.25.
Claude Haiku~4.5 comes in second at 60.78\%, despite having the
worst fertility of any instruction tuned model I tested (7.73).
Anon-Indic-GPT2 has the best fertility of all, 1.60, and still
scores near chance at 26.97\%, which fits its profile as a 110M
parameter base language model with no instruction tuning.

This decoupling matters for how the field should prioritize its
work. A better tokenizer will not make a model understand Telugu
better, and a strong Telugu benchmark score does not mean Telugu
speakers are paying a fair price. Tokenizer fairness has to be its
own engineering goal: include Telugu and other Indic subword units
in the vocabulary at pretraining time, rather than hoping it falls
out of scaling up the model.

Sarvam and Claude make an instructive comparison here. Sarvam-2B has
a favorable script fairness gap, 0.53, but its zero shot MCQ
accuracy was unusable; only 7.9\% of responses were even parseable,
so I excluded it. The tokenizer is well calibrated; the model just
was not built to follow instructions. A future model that paired
Sarvam's tokenizer with instruction tuning at frontier scale would
fix both problems at once.

\subsection{Limitations}
\label{sec:limitations}

\paragraph{Single romanization scheme.}
My romanized corpus relies only on ITRANS transliteration. Other
romanization conventions, ISO~15919, or the informal phonetic
spellings common on social media, might produce different fertility
numbers. I treat ITRANS as a reasonable worst case for general
purpose tokenizers, but the exact size of the gap will shift with
romanization style.

\paragraph{Corpus size and domain coverage.}
My corpus of 2,789 sentences is large enough for significance
testing but still small next to production scale corpora. The
native formal register comes only from Wikipedia, which skews
toward formal vocabulary and may not capture colloquial
morphological forms well. A larger corpus spanning news, social
media, and literature would give more precise fertility estimates.

\paragraph{Single language.}
I only study Telugu here. The script fairness gap will probably
look different for other Dravidian languages such as Tamil, Kannada,
and Malayalam, since they differ in script complexity, morphological
richness, and how much they show up in tokenizer training corpora.
Extending this work across languages is left for the future.

\paragraph{Sarvam-2B evaluation.}
The exclusion of Sarvam-2B from the accuracy analysis---necessitated
by the fact that 92.1\% of its responses were unparseable---limits
what I can say about whether Indic optimised tokenizers paired with
instruction tuned models close the capability gap. Future work should score Sarvam models
with the log likelihood method instead of generative MCQ.

\paragraph{Minimal pair validation.}
The 62 minimal pairs were built using Telugu morphological grammars
and checked by one native speaker. A larger validation effort across
more speakers and dialects would make the morphological analysis
stronger.


\section{Conclusion}
\label{sec:conclusion}
This paper presented a systematic empirical audit of tokenizer
fertility and script fairness for Telugu, covering six LLM
tokenizers and three corpus registers. My findings answer all
four research questions in turn.

Fertility varies by up to 13.8$\times$ across tokenizers on the same
native Telugu text: Anon-Indic-GPT2 sits at one extreme (1.60
tokens/word) and GPT-2 at the other (22.08 tokens/word) (RQ1). The
script fairness gap is statistically significant, and which way it
runs depends on the tokenizer's origin. General purpose Western
tokenizers impose a statistically significant native script penalty
between 1.56$\times$ and 4.66$\times$ ($p < 10^{-165}$,
$n = 1{,}000$). Tokenizers trained on Indic data invert that bias:
native script ends up cheaper than romanized input ($\mathcal{G}$
between 0.29 and 0.68), confirmed by Wilcoxon tests in the opposing
direction ($p < 10^{-165}$, $n = 1{,}000$) (RQ2). The penalty also follows morphological
structure: honorific suffixes and verb agglutination chains take
the highest token counts of any category, the same constructions
that define Telugu as an agglutinative language (RQ3). Tokenizer
fertility and downstream MILU accuracy are decoupled: the model
with the best tokenizer scores near chance, 26.97\%, while the
model with the worst tokenizer among instruction tuned models
scores 60.78\%. That gap is the clearest evidence in the paper that
the token tax is about pricing, not capability (RQ4).

These findings point to two concrete recommendations for tokenizer
designers. First, Telugu and other Dravidian subword units belong
in the vocabulary itself, not left for byte level fallback to
handle. Sarvam shows this is possible at a 64k vocabulary size
without giving up cross lingual coverage across 22 Indic languages.
Second, the script fairness gap should become a standard metric
reported alongside fertility for multilingual tokenizers
\citep{rust-etal-2021-how, petrov-etal-2023-language}, since the two
numbers can diverge substantially, as my results show.

The token tax comes down to a structural inequity in how LLM access
is priced and scoped. Telugu has 96 million speakers, most of whom
access AI services on limited budgets, and they pay more per query,
get less context per rupee, and remain underrepresented in both
tokenizer training data and model pretraining corpora
\citep{census2011, doddapaneni-etal-2023-towards}. Closing that gap
takes deliberate engineering choices at the tokenizer level, choices
that do not depend on model scale at all. I release my corpus,
minimal pair dataset, evaluation scripts, and every other
experimental artifact, so other researchers can build on this work
for Telugu and for other morphologically rich, low resource
languages.

% Bibliography entries for the entire Anthology, followed by custom entries
%\bibliography{anthology,custom}
% Custom bibliography entries only
\bibliography{custom}

\appendix


\end{document}
